\section{引言 Introduction}
\label{sec:intro}

多语言预训练模型的词表通常由高资源语言的语料统计得到。当这一词表被应用到
形态丰富但语料稀缺的语言时，单词往往被切分成大量无语义的片段。我们在四种
低资源语言上测得，平均每个词被切分为 4.7 个词元，而同一模型在英语上仅为
1.3 个。碎片化直接抬高了序列长度，也削弱了模型对词内结构的建模能力。

A tokenizer inherited from high-resource data fragments morphologically rich
languages. The fragmentation is not merely an efficiency problem. Longer
sequences consume the context budget, and the model must rediscover word
boundaries that the tokenizer has destroyed.

本文的出发点是：形态规则在亲缘语言之间是可迁移的，即使词汇本身不可迁移。
我们据此构造一个共享形态先验，并用它约束目标语言的词表学习过程。

\subsection{贡献 Contributions}

\begin{itemize}
  \item 提出一种形态先验约束下的分词迁移算法，仅需十万句目标语言语料。
  \item 在四种低资源语言上验证该方法，平均词元长度提高百分之三十一。
  \item 明确给出方法失效的两种情形，见第~\ref{sec:experiments} 节。
\end{itemize}
